\section{Discussion and Open Problems}
\label{sec:future}

Non-structural subtype entailment with one constructor of any arity and
variance is decidable over arbitrary, regular, and finite trees. The proof
reduces non-entailment to holes of finite readers, uses the rejected-path
property of these readers to bound the rejected tail and hence the depth of
a block hierarchy, and then shortens all blocks of each level
simultaneously. The questions below concern the cost of this procedure, the
role of the rejected-path property, and richer signatures and conclusions.

\subsection{Complexity and Sufficient Tests}
\Cref{sec:complexity} places entailment in 2-EXPSPACE for every arity and
variance, in each domain, and for a covariant constructor of each fixed arity
$n\ge2$ it is PSPACE-hard (\cref{thm:cx-hard}). Three questions remain open.
First, whether entailment for $n\ge2$ is in PSPACE, or at least below
2-EXPSPACE. Second, whether the PSPACE bound at arity one
(\cref{app:small-arity}) is tight. Third, whether a left or right label
reader that has a hole always has one of length singly exponential in its
state count $N$. The long holes of \cref{prop:cx-length} do not settle the
third question: their length $2^{\Omega(k)}$ is only polynomial in
$N=2^{4k+1}$, whereas the length bound $H$ is triply exponential in $k$.
In the fully covariant case, self events vanish, but the dictionary bound
does not exploit this to improve the estimate. Reducing the depth,
sharpening the dictionary bound, or using the label readers' stronger
restrictions on admission (\cref{sec:interface}) may narrow the gap between
PSPACE and 2-EXPSPACE.

The sufficient test of \cref{lem:branching} could also serve as an early
non-entailment test in a simplifier; its cost is polynomial in the explicit
reader graph, which can itself be exponential in the input.

\subsection{The General Hole Problem}\label{sec:general-holes}
For arbitrary readers, the remaining question can be stated precisely:
is there a computable function $B(N)$ such that every $N$-state reader over
any finite alphabet with a hole has one with rejected tail at most $B(N)$?
Under the rejected-path property, $B(N)=4N-2$ works (\cref{thm:brt}).
The depth and compression theorems hold for every reader, so such a
function would decide hole existence by searching to
$\holeBound(N,2N(B(N)+1)-1)$. Conversely, a decision procedure for hole
existence would give such a function, even uniformly over finite alphabets.

The uniformity rests on an alphabet reduction. There are at most $N^N$
letter transitions $Q\to Q$, and for hole existence it suffices to retain
at most two letters with each transition. Indeed, let $w$ be a hole. If no
letter of $w$ shares its transition with another letter of the alphabet,
$w$ already uses only retained letters. Otherwise, consider the last
position whose letter shares its transition. Replace every letter by the
first retained letter with the same transition, except at that position,
where a second letter with that transition serves as a marker occurring
nowhere else. The run is unchanged. The marker excludes every nonterminal
comparison ending at or before the marked position. Every comparison ending
after it compares segments made of letters with unique transitions, which
the replacement leaves unchanged, so it already existed in $w$ with the same
states. The new word is therefore a hole with the same rejected tail.
Conversely, injecting the retained alphabet into the original one preserves
holes. Up to renaming, there are thus finitely many $N$-state readers to
consider, each with at most $2N^N$ letters. Given a hole-existence decision
procedure, decide which of them have holes, enumerate a hole for each
positive answer, and take the maximum rejected-tail length, or zero if
there are none. This maximum is a computable bound for every $N$-state
reader. Thus decidability of hole existence for all finite readers and the
computable-tail question are equivalent. The linear bound of \cref{thm:brt}
answers the question only under the rejected-path property.

Hole languages need not be regular, even for label readers. Take the
binary covariant constructor, the single literal $f(y,x)\le y$, and the
query $x\le y$ from \cref{sec:intro}. A word is right unsafe exactly when it
has the form $0^i1z$ with $z\not\pre(0^i1)^\omega$. Indeed, let $\sigma$ be a
solution in which $\sigma y$ has a bottom prefix on $0^i1z$. Since $\sigma y$
lies below its $0$-child, its nodes on $0^i$ are constructors, each lying
below its own $0$-child, so its subtree at $0^i1$ lies above $\sigma x$. The
bottom prefix lies at $0^i1u$ with $u\pre z$, and then $\sigma x$ has a
bottom prefix on $u$. If $z\pre(0^i1)^\omega$, then $u$ and $0^i1z$ prefix
the same infinite word, so $u\pre 0^i1z$ and the word is not right unsafe.
Conversely, let $\sigma x$ have constructors exactly at the prefixes of
$w=0^i1z$ and bottom leaves elsewhere, and let $\sigma y=f(\sigma y,\sigma
x)$. Then $\sigma y$ has a constructor at $w$ iff $z\pre 0^i1z$, that is,
iff $z\pre(0^i1)^\omega$; otherwise it has a bottom prefix on $w$. The
system is satisfiable, so by \cref{thm:holes} these words are exactly the
holes of the right label reader
$(\mathcal{M}_{\mathsf r},\Adm_{\mathsf r},\Rej_{\mathsf r})$. Their
intersection with $0^*10^*$ is $\{0^i10^j: j>i\}$, which is not regular.
Our finite search decides existence under the rejected-path property
without constructing an automaton for the whole hole language.

\subsection{Several Constructors}
With several constructor labels, the label order need no longer be the
chain $\Bot<f<\Top$: different constructors may be incomparable.
The label-clash analysis must then distinguish incompatible constructors,
and a least-shape construction must choose a compatible constructor as
well as decide where a constructor is needed. Path polarity depends on
which constructor occurs at each node, and available child positions depend
on that constructor's arity. The normalization and path decompositions
would have to track this additional information.

Two parts of the present proof may transfer more readily. Signed
translation with polarity selection is a general device for variance:
solve a covariant system on two copies of each variable, then select a
solution satisfying the required relation between the copies. The
alphabet-generic word bounds remain available for any semantic reduction
that produces finite readers with the rejected-path property, for instance
from string constraints or automata with equality constraints that yield
suffix-periodicity conditions. The unresolved part is obtaining such a
reduction. Several nullary base types together with one binary constructor
form a natural next case for type inference; no result for that signature
is claimed here.

\subsection{Richer Entailment Fragments}
\label{sec:conclusion}

The judgment studied here has one variable inequality as its conclusion;
model-theoretically, it is the validity of a universal implication
$\forall\vec x\,(\phi\Rightarrow x\le y)$. Natural extensions include
existential entailment $\phi\models\exists\vec z.\,\psi$ and entailment of a
disjunction of inequalities. These conclusions ask either for suitable
values of $\vec z$ in each solution or for one of several inequalities to
hold; neither is decided by testing the component inequalities
independently. Undecidability of the full first-order theory
\citep{SuAikenNiehrenPriesnitzTreinen02} does not settle these individual
fragments. Determining which such fragments admit finite reader reductions
remains an open direction.

